\documentclass[preview]{standalone}

\usepackage{amsmath}
\usepackage{amssymb}
\usepackage{stellar}
\usepackage{definitions}
\usepackage{bettelini}

\begin{document}

\id{probability-convergence}
\genpage

\section{Modes of convergence}

\begin{snippetdefinition}{convergence-in-distribution-definition}{Convergence in distribution}
    Let \((X_n)_{n\geq 1}\) be a sequence of real-valued
    \randomvariable[random variables] and let \(X\) be a real-valued
    \randomvariable, all defined on possibly different \probabilityspace[probability spaces].
    Let \(F_n\) be the \distributionfunction of \(X_n\) and let \(F\) be the
    \distributionfunction of \(X\). We say that \(X_n\)
    \emph{converges in distribution} to \(X\), and we write
    \[
        X_n \xrightarrow{d} X,
    \]
    if
    \[
        \lim_{n\to\infty}F_n(x)=F(x)
    \]
    for every continuity point \(x\) of \(F\).
\end{snippetdefinition}

\begin{snippetdefinition}{convergence-in-probability-definition}{Convergence in probability}
    Let \((\Omega,\mathcal{F},\mathbb{P})\) be a \probabilityspace.
    Let \((X_n)_{n\geq 1}\) be a sequence of real-valued
    \randomvariable[random variables] on \(\Omega\), and let \(X\) be a
    real-valued \randomvariable on \(\Omega\). We say that \(X_n\)
    \emph{converges in probability} to \(X\), and we write
    \[
        X_n \xrightarrow{\mathbb{P}} X,
    \]
    if, for every \(\varepsilon>0\),
    \[
        \lim_{n\to\infty}\mathbb{P}\left(|X_n-X|>\varepsilon\right)=0.
    \]
\end{snippetdefinition}

\begin{snippetdefinition}{almost-sure-convergence-definition}{Almost sure convergence}
    Let \((\Omega,\mathcal{F},\mathbb{P})\) be a \probabilityspace.
    Let \((X_n)_{n\geq 1}\) be a sequence of real-valued
    \randomvariable[random variables] on \(\Omega\), and let \(X\) be a
    real-valued \randomvariable on \(\Omega\). We say that \(X_n\)
    \emph{converges almost surely} to \(X\), and we write
    \[
        X_n \xrightarrow{\mathrm{a.s.}} X,
    \]
    if there exists an \event \(\Omega_0\in\mathcal{F}\) such that
    \(\mathbb{P}(\Omega_0)=1\) and
    \[
        \lim_{n\to\infty}X_n(\omega)=X(\omega)
    \]
    for every \(\omega\in\Omega_0\).
\end{snippetdefinition}

\section{Relations between modes}

\begin{snippettheorem}{convergence-in-probability-implies-distribution}{Convergence in probability implies convergence in distribution}
    Let \((\Omega,\mathcal{F},\mathbb{P})\) be a \probabilityspace.
    Let \((X_n)_{n\geq 1}\) be a sequence of real-valued
    \randomvariable[random variables] on \(\Omega\), and let \(X\) be a
    real-valued \randomvariable on \(\Omega\). If
    \[
        X_n\xrightarrow{\mathbb{P}}X,
    \]
    then
    \[
        X_n\xrightarrow{d}X.
    \]
\end{snippettheorem}

\begin{snippetproof}{convergence-in-probability-implies-distribution-proof}{convergence-in-probability-implies-distribution}{Convergence in probability implies convergence in distribution}
    Let \(F_n\) be the \distributionfunction of \(X_n\), and let \(F\) be
    the \distributionfunction of \(X\). Fix a continuity point \(x\) of
    \(F\), and fix \(\varepsilon>0\). Since
    \[
        \{X_n\leq x\}
        \subseteq
        \{X\leq x+\varepsilon\}\union\{|X_n-X|>\varepsilon\},
    \]
    we have
    \[
        F_n(x)
        \leq F(x+\varepsilon)
        +\mathbb{P}\left(|X_n-X|>\varepsilon\right).
    \]
    Taking the upper limit as \(n\to\infty\) gives
    \[
        \limsup_{n\to\infty}F_n(x)\leq F(x+\varepsilon).
    \]
    Similarly,
    \[
        \{X\leq x-\varepsilon\}
        \subseteq
        \{X_n\leq x\}\union\{|X_n-X|>\varepsilon\},
    \]
    and hence
    \[
        F(x-\varepsilon)
        \leq
        F_n(x)+\mathbb{P}\left(|X_n-X|>\varepsilon\right).
    \]
    Taking the lower limit as \(n\to\infty\) gives
    \[
        F(x-\varepsilon)\leq\liminf_{n\to\infty}F_n(x).
    \]
    Therefore
    \[
        F(x-\varepsilon)
        \leq
        \liminf_{n\to\infty}F_n(x)
        \leq
        \limsup_{n\to\infty}F_n(x)
        \leq
        F(x+\varepsilon).
    \]
    Since \(F\) is continuous at \(x\), letting \(\varepsilon\to0^+\)
    yields
    \[
        \lim_{n\to\infty}F_n(x)=F(x).
    \]
    Thus \(X_n\xrightarrow{d}X\).
\end{snippetproof}

\begin{snippetexample}{distribution-convergence-not-probability-example}{Convergence in distribution need not imply convergence in probability}
    Let \(X\) be a \randomvariable with
    \(X\sim\operatorname{Bernoulli}(1/2)\), let \(Y=1-X\), and define
    \(X_n=X\) for every \(n\geq1\). Then \(X_n\eqdist Y\) for every
    \(n\geq1\), because both \(X_n\) and \(Y\) have the
    \bernoullidistribution with parameter \(1/2\). Hence
    \[
        X_n\xrightarrow{d}Y.
    \]
    However, for every \(0<\varepsilon<1\),
    \[
        \mathbb{P}\left(|X_n-Y|>\varepsilon\right)
        =
        \mathbb{P}\left(|2X-1|>\varepsilon\right)=1.
    \]
    Thus \(X_n\) does not converge in probability to \(Y\).
\end{snippetexample}

\section{Event limits}

\begin{snippetdefinition}{limsup-liminf-events-definition}{Limit superior and limit inferior of events}
    Let \((\Omega,\mathcal{F},\mathbb{P})\) be a \probabilityspace, and
    let \((A_n)_{n\geq 1}\) be a sequence of \event[events]. The
    \emph{limit superior} and \emph{limit inferior} of \((A_n)\) are
    \[
        \limsup_{n\to\infty}A_n
        =
        \bigcap_{m=1}^{+\infty}\bigcup_{n=m}^{+\infty}A_n
    \]
    and
    \[
        \liminf_{n\to\infty}A_n
        =
        \bigcup_{m=1}^{+\infty}\bigcap_{n=m}^{+\infty}A_n.
    \]
    The \event \(\limsup_{n\to\infty}A_n\) occurs when infinitely many
    \(A_n\) occur. The \event \(\liminf_{n\to\infty}A_n\) occurs when
    all \(A_n\) occur from some index onward.
\end{snippetdefinition}

\begin{snippettheorem}{event-fatou-inequalities}{Fatou inequalities for events}
    Let \((\Omega,\mathcal{F},\mathbb{P})\) be a \probabilityspace, and
    let \((A_n)_{n\geq1}\) be a sequence of \event[events]. Then
    \[
        \mathbb{P}\left(\liminf_{n\to\infty}A_n\right)
        \leq
        \liminf_{n\to\infty}\mathbb{P}(A_n)
        \leq
        \limsup_{n\to\infty}\mathbb{P}(A_n)
        \leq
        \mathbb{P}\left(\limsup_{n\to\infty}A_n\right).
    \]
    In particular, if
    \[
        \liminf_{n\to\infty}A_n=\limsup_{n\to\infty}A_n=A,
    \]
    then
    \[
        \lim_{n\to\infty}\mathbb{P}(A_n)=\mathbb{P}(A).
    \]
\end{snippettheorem}

\begin{snippetproof}{event-fatou-inequalities-proof}{event-fatou-inequalities}{Fatou inequalities for events}
    Put
    \[
        S_m=\bigcup_{n=m}^{+\infty}A_n
        \quad\text{and}\quad
        L_m=\bigcap_{n=m}^{+\infty}A_n.
    \]
    Then \((S_m)\) is decreasing, \((L_m)\) is increasing, and
    \[
        \limsup_{n\to\infty}A_n=\bigcap_{m=1}^{+\infty}S_m,
        \quad
        \liminf_{n\to\infty}A_n=\bigcup_{m=1}^{+\infty}L_m.
    \]
    Since \(L_m\subseteq A_n\subseteq S_m\) for every \(n\geq m\),
    \[
        \mathbb{P}(L_m)
        \leq
        \inf_{n\geq m}\mathbb{P}(A_n)
        \leq
        \sup_{n\geq m}\mathbb{P}(A_n)
        \leq
        \mathbb{P}(S_m).
    \]
    Letting \(m\to\infty\), and using continuity of probability from
    below for \((L_m)\) and from above for \((S_m)\), gives the stated
    inequalities.
\end{snippetproof}

\begin{snippettheorem}{borel-cantelli-lemma}{Borel-Cantelli lemma}
    Let \((\Omega,\mathcal{F},\mathbb{P})\) be a \probabilityspace, and
    let \((A_n)_{n\geq1}\) be a sequence of \event[events]. If
    \[
        \sum_{n=1}^{+\infty}\mathbb{P}(A_n)<+\infty,
    \]
    then
    \[
        \mathbb{P}\left(\limsup_{n\to\infty}A_n\right)=0.
    \]
\end{snippettheorem}

\begin{snippetproof}{borel-cantelli-lemma-proof}{borel-cantelli-lemma}{Borel-Cantelli lemma}
    For every \(m\geq1\),
    \[
        \limsup_{n\to\infty}A_n
        \subseteq
        \bigcup_{n=m}^{+\infty}A_n.
    \]
    Hence, by countable subadditivity,
    \[
        \mathbb{P}\left(\limsup_{n\to\infty}A_n\right)
        \leq
        \sum_{n=m}^{+\infty}\mathbb{P}(A_n).
    \]
    Since the series \(\sum_{n=1}^{+\infty}\mathbb{P}(A_n)\) converges,
    its tails converge to \(0\). Letting \(m\to\infty\) gives
    \[
        \mathbb{P}\left(\limsup_{n\to\infty}A_n\right)=0.
    \]
\end{snippetproof}

\section{Almost sure criteria}

\begin{snippettheorem}{almost-sure-convergence-limsup-criterion}{Limsup criterion for almost sure convergence}
    Let \((\Omega,\mathcal{F},\mathbb{P})\) be a \probabilityspace.
    Let \((X_n)_{n\geq1}\) be a sequence of real-valued
    \randomvariable[random variables] on \(\Omega\), and let \(X\) be a
    real-valued \randomvariable on \(\Omega\). For every
    \(\varepsilon>0\), define
    \[
        A_n(\varepsilon)=\{|X_n-X|>\varepsilon\}.
    \]
    Then
    \[
        X_n\xrightarrow{\mathrm{a.s.}}X
    \]
    if and only if, for every \(\varepsilon>0\),
    \[
        \mathbb{P}\left(\limsup_{n\to\infty}A_n(\varepsilon)\right)=0.
    \]
\end{snippettheorem}

\begin{snippetproof}{almost-sure-convergence-limsup-criterion-proof}{almost-sure-convergence-limsup-criterion}{Limsup criterion for almost sure convergence}
    For \(\omega\in\Omega\), the sequence \(X_n(\omega)\) converges to
    \(X(\omega)\) if and only if, for every rational
    \(\varepsilon>0\), the inequality
    \(|X_n(\omega)-X(\omega)|>\varepsilon\) holds only finitely many
    times. Hence
    \[
        \{X_n\not\to X\}
        =
        \bigcup_{\substack{\varepsilon\in\rationalnumbers\\ \varepsilon>0}}
        \limsup_{n\to\infty}A_n(\varepsilon).
    \]
    Therefore \(\mathbb{P}(\{X_n\not\to X\})=0\) if and only if
    every \(\limsup_{n\to\infty}A_n(\varepsilon)\), with
    \(\varepsilon\in\rationalnumbers\) and \(\varepsilon>0\), has
    probability \(0\). This is equivalent to the same condition for
    every real \(\varepsilon>0\).
\end{snippetproof}

\begin{snippetcorollary}{summability-implies-almost-sure-convergence}{Summability criterion for almost sure convergence}
    Let \((\Omega,\mathcal{F},\mathbb{P})\) be a \probabilityspace.
    Let \((X_n)_{n\geq1}\) be a sequence of real-valued
    \randomvariable[random variables] on \(\Omega\), and let \(X\) be a
    real-valued \randomvariable on \(\Omega\). If, for every
    \(\varepsilon>0\),
    \[
        \sum_{n=1}^{+\infty}
        \mathbb{P}\left(|X_n-X|>\varepsilon\right)
        <+\infty,
    \]
    then
    \[
        X_n\xrightarrow{\mathrm{a.s.}}X.
    \]
\end{snippetcorollary}

\begin{snippetproof}{summability-implies-almost-sure-convergence-proof}{summability-implies-almost-sure-convergence}{Summability criterion for almost sure convergence}
    Fix \(\varepsilon>0\) and put
    \[
        A_n(\varepsilon)=\{|X_n-X|>\varepsilon\}.
    \]
    By hypothesis,
    \[
        \sum_{n=1}^{+\infty}\mathbb{P}(A_n(\varepsilon))<+\infty.
    \]
    The \borelcantellilemma gives
    \[
        \mathbb{P}\left(\limsup_{n\to\infty}A_n(\varepsilon)\right)=0.
    \]
    Since this holds for every \(\varepsilon>0\), the
    \almostsureconvergence criterion gives
    \(X_n\xrightarrow{\mathrm{a.s.}}X\).
\end{snippetproof}

\begin{snippettheorem}{almost-sure-convergence-implies-probability}{Almost sure convergence implies convergence in probability}
    Let \((\Omega,\mathcal{F},\mathbb{P})\) be a \probabilityspace.
    Let \((X_n)_{n\geq1}\) be a sequence of real-valued
    \randomvariable[random variables] on \(\Omega\), and let \(X\) be a
    real-valued \randomvariable on \(\Omega\). If
    \[
        X_n\xrightarrow{\mathrm{a.s.}}X,
    \]
    then
    \[
        X_n\xrightarrow{\mathbb{P}}X.
    \]
\end{snippettheorem}

\begin{snippetproof}{almost-sure-convergence-implies-probability-proof}{almost-sure-convergence-implies-probability}{Almost sure convergence implies convergence in probability}
    Fix \(\varepsilon>0\), and put
    \[
        A_n=\{|X_n-X|>\varepsilon\}.
    \]
    By the \almostsureconvergence criterion,
    \[
        \mathbb{P}\left(\limsup_{n\to\infty}A_n\right)=0.
    \]
    Since
    \[
        \limsup_{n\to\infty}\mathbb{P}(A_n)
        \leq
        \mathbb{P}\left(\limsup_{n\to\infty}A_n\right),
    \]
    we get
    \[
        \lim_{n\to\infty}\mathbb{P}\left(|X_n-X|>\varepsilon\right)=0.
    \]
    Thus \(X_n\xrightarrow{\mathbb{P}}X\).
\end{snippetproof}

\begin{snippetexample}{probability-convergence-not-almost-sure-example}{Convergence in probability need not imply almost sure convergence}
    Let \((X_n)_{n\geq1}\) be independent \randomvariable[random variables]
    such that
    \[
        X_n\sim\operatorname{Bernoulli}(1/n).
    \]
    Then \(X_n\xrightarrow{\mathbb{P}}0\). Indeed, if
    \(0<\varepsilon<1\), then
    \[
        \mathbb{P}(|X_n|>\varepsilon)=\mathbb{P}(X_n=1)=\frac{1}{n},
    \]
    while for \(\varepsilon\geq1\) this probability is \(0\).

    However \(X_n\) does not converge almost surely to \(0\). Let
    \(A_n=\{X_n=1\}\). For \(m\geq2\),
    \[
        \mathbb{P}\left(\bigcap_{n=m}^{k} A_n^c\right)
        =
        \prod_{n=m}^{k}\left(1-\frac{1}{n}\right)
        =
        \frac{m-1}{k}.
    \]
    Letting \(k\to\infty\), we get
    \[
        \mathbb{P}\left(\bigcap_{n=m}^{+\infty} A_n^c\right)=0.
    \]
    Hence
    \[
        \mathbb{P}\left(\bigcup_{n=m}^{+\infty}A_n\right)=1
    \]
    for every \(m\geq2\), and therefore
    \[
        \mathbb{P}\left(\limsup_{n\to\infty}A_n\right)=1.
    \]
    Since \(A_n=\{|X_n-0|>1/2\}\), the
    \almostsureconvergence criterion shows that
    \(X_n\not\xrightarrow{\mathrm{a.s.}}0\).
\end{snippetexample}

\end{document}
